\section{Design}
\label{sec:design}

\sysname{} has two halves: a compile-time analysis that produces structure, and a
runtime that spends idle GPU time against that structure.

\subsection{Static analysis over byLLM programs}

The analysis runs as a compiler pass and requires no annotation from the
developer and no trained model. It produces three artifacts.

\paragraph{Call topology.} The pass deduces the topology of byLLM calls: which
call sites can reach which other call sites, expanded as a BFS-layered reachable
set rather than a single next-step guess. The result is a may-happen relation, so
a branch whose outcome is not statically decidable contributes all of its
successors rather than none of them.

\paragraph{Argument provenance.} For every parameter of every call site, the pass
records where the value that will be bound to it is produced. This is the
artifact that distinguishes \sysname{} from next-agent prediction: knowing that
call site $v$ will run is far less actionable than knowing that the value being
computed right now is the third argument of $v$, and is therefore prefillable
right now.

\paragraph{Invariant/variant split.} Each runtime prompt is partitioned into the
portion that is fixed at compile time and the portion that depends on runtime
values. The invariant portion is the part that can be materialized without any
speculation at all.

\subsection{Cache-aware prompt layout}

Because the compiler owns prompt assembly, it does not have to accept the prompt
as given. byLLM's stock rendering interleaves invariant and variant content, which
fragments what would otherwise be a long shared prefix: a single early-positioned
dynamic value invalidates every token after it for reuse purposes. The pass
rewrites the layout so that zones are ordered by expected stability --- invariant
instruction and schema content first, graph context next, and a trailing binding
zone whose internal order follows the arrival order recovered by the provenance
analysis. The effect is that the prefix that is reusable is also contiguous, which
is what a prefix cache is able to exploit.

This is an authority a serving-layer system structurally does not have. A system
that sits between the agent framework and the
engine~\cite{zhang2026cachescout,yu2026pythia} can decide what to keep, evict, or
warm, but it receives the prompt already rendered.

\subsection{Runtime: prefilling in engine idle windows}

The runtime consumes the compiler's artifacts and issues prefill work when the
engine would otherwise be underutilized. It targets two kinds of idle. \emph{Temporal}
idle is the wall-clock gap in which no work is outstanding for the request,
opened by a tool call or a retrieval round trip. \emph{Spatial} idle is capacity
within a scheduling step that is not consumed by the current batch; \sysname{}
reaches it through chunked prefill~\cite{agrawal2024sarathi}, which lets prefill
tokens be admitted incrementally alongside ongoing work rather than as an
indivisible unit.

Within an idle window, the runtime builds the longest knowable prefix of a
reachable future call: its invariant zone, the graph context that is already
determined, the argument values that provenance says are already bound, and the
grammar compilation required for constrained decoding of the call's structured
output~\cite{dong2025xgrammar}. This is a strictly larger unit than a fixed
anchor. Where a warmup-based approach materializes the recurring context of a
predicted agent, \sysname{} additionally materializes the speculated argument
content, which is the part of the prompt that content-level approaches set aside.

\subsection{Proactive binding of early-computed arguments}

The single most productive case falls directly out of provenance: an argument
that is prepared \emph{before} the byLLM function that consumes it is called.
Such a value is not speculative at all --- it exists, it is final, and its
consumer is statically known. The runtime prefills it as soon as it is produced,
so by the time control reaches the consuming call site the corresponding KV is
already resident. Section~\ref{sec:eval} shows that this case accounts for the
bulk of the measured gain: call sites whose prompts accumulate the outputs of
earlier stages are precisely the ones that improve most.

\subsection{Resolving supervisor branches}

A supervisor that selects one of several subagents creates a branch that the
static topology can only over-approximate. \sysname{} includes a training-free
mechanism to narrow it: truncate the supervisor's reasoning and fork a probe
decode to see which successor it is heading toward, or equivalently run a draft
model to speculate the routing decision, and prefill that successor first. The
mechanism requires no additional training and no workload-specific tuning.

We note explicitly that this component is \emph{not enabled} in the measurements
reported in Section~\ref{sec:eval}: the configuration evaluated there uses
compiler-derived speculative prefill only, with no probe and no deploy-time
prewarm. The consequence is visible in the results, where the call sites
immediately downstream of the routing decision show no improvement.

\subsection{Correctness model}

Speculation in \sysname{} is advisory and byte-verified. The prompt that is
served is byLLM's own rendering of the call, unchanged; the speculatively
constructed KV is used only if the tokens it covers match that rendering exactly.
A mispredicted branch, a stale value, or a layout mismatch therefore degrades to
a cache miss and a normal cold prefill. There is no configuration in which
speculation changes the prompt the model sees, and hence none in which it changes
the program's output.
